\documentclass[11pt]{article}
\usepackage[T1]{fontenc}
\usepackage{lmodern}
\usepackage{amsmath,amssymb,mathtools}
\usepackage[a4paper,margin=28mm]{geometry}
\usepackage{microtype}
\usepackage{xcolor}
\usepackage{enumitem}
\usepackage{listings}
\usepackage[hidelinks]{hyperref}
\usepackage[nameinlink,noabbrev]{cleveref}
\setlength{\emergencystretch}{3em}
\definecolor{codegray}{RGB}{245,245,245}
\newcommand{\workflowstep}[1]{\subsection{Workflow Step #1}}
\newcommand{\taupar}{\tau}
\newcommand{\taumod}{\tau_{\mathrm{mod}}}
\newcommand{\EJ}{E_J}
\lstdefinestyle{parigp}{
  language={},
  basicstyle=\ttfamily\small,
  backgroundcolor=\color{codegray},
  frame=single,
  breaklines=true,
  columns=fullflexible,
  keepspaces=true,
  showstringspaces=false
}
\title{From Direct Algebraic Generating Functions to Elliptic Sigma and
Theta Representations of Somos--4 Sequences\\[0.4em]
\large A Constructive Workflow}
\author{Thomas Scheuerle}
\date{}

\begin{document}
\maketitle

\begin{abstract}
A direct five-parameter algebraic generating function can be constructed
whose Hankel transform follows a prescribed nondegenerate Somos--4 orbit.
After removal of a square factor, its discriminant supplies the genus-one
curve
\[
  \EJ:\qquad y^2=(X^2-JX+\taupar)^2-4X,
\]
which is subsequently identified with the invariant curve of the reduced
Somos--4 QRT dynamics.
The elliptic-function theory behind Somos--4 sequences is classical; the
practical difficulty lies in moving from concrete recurrence data to a
numerical sigma- or theta-function representation without losing track of
normalizations. The relevant ingredients are scattered across several
sources and conventions. This paper brings them together into one
constructive and reproducible workflow:
\[
 (\lambda,m,r,\eta,\taupar)
 \longmapsto J
 \longmapsto
 \begin{cases}
   \text{discriminant quartic},\\
   \text{QRT invariant curve}
 \end{cases}
 \longmapsto \EJ
 \longmapsto P
 \longmapsto (\omega_1,\omega_2;z_0,\delta)
 \longmapsto \sigma
 \longmapsto \theta_1.
\]
The emphasis is on normalization choices, computational checks, and an
executable PARI/GP implementation. No new proof of the classical
sigma-function representation is intended; the contribution is a coherent
end-to-end procedure that starts from the algebraic generating function and
ends with a numerically verified sigma description and an explicit conversion
to the Jacobi theta function $\theta_1$.
\end{abstract}

\section{Introduction}
The relation between Somos--4 recurrences, elliptic curves, and elliptic
functions is classical; see, in particular, \cite{Ward1948,Hone2005,
VanderPoortenSwart2006,HoneSwart2008}. Generic Somos--4 orbits admit
representations in terms of the Weierstrass sigma function
\cite{Hone2005}. Nevertheless, a reader beginning with a concrete sequence
still faces a substantial practical problem: the relevant constructions
occur in different normalizations, coordinate systems, and parametrizations,
and the intermediate steps needed for a numerical implementation are rarely
presented together.

The purpose of this paper is not to prove a new elliptic-function theorem.
Rather, it is to make the established theory directly usable. Our starting
point is the direct algebraic generating function constructed in
\cite{Scheuerle2026}. Its discriminant naturally produces a genus-one curve
associated with the Somos--4 data. From there we describe how to obtain a
Weierstrass model, compute the period lattice, identify the translation point,
determine the Abel--Jacobi parameters, and finally write explicit sigma- and
representations in terms of the Weierstrass sigma function and the Jacobi
theta function $\theta_1$. The continued-fraction and Hankel viewpoint of
\cite{Hone2020} is complementary: it explains structurally why genus-one
continued-fraction data, Hankel determinants, and Somos--4 dynamics are
governed by translation on the same Jacobian.

The paper is organized as a mathematical workflow rather than a
theorem--proof sequence. Classical facts are stated with references and are
not reproved. The main concern is to keep the conventions aligned and to
expose the checks that catch a wrong branch, period, sign, or scaling.
PARI/GP fragments accompany the construction. The appendix contains an
end-to-end verification of the elliptic and sigma-function stages beginning
with the reduced parameters $(J,\taupar)$. It does not duplicate the original
generating-function discriminant calculation, and the theta conversion is
derived explicitly but not evaluated independently by the supplied script.

\subsection{Scope and contribution}
This paper provides:
\begin{itemize}[leftmargin=2em]
  \item one route from the five input parameters to a genus-one curve and its
        identification with the invariant curve of the reduced QRT dynamics;
  \item an explicit record of the normalization choices needed to pass from
        a quartic model to the Weierstrass sigma function and the Jacobi theta
        function $\theta_1$;
  \item computational checkpoints at each implemented stage;
  \item worked examples comparing recurrence and sigma-function evaluations,
        together with a fully explicit algebraic theta conversion; and
  \item an end-to-end PARI/GP workflow, beginning with $(J,\taupar)$, suitable
        for adaptation.
\end{itemize}
It does not claim a new proof of the classical Somos--4 sigma formula, a new
integrability result, or a new theorem on elliptic curves.

\section{Input data and the direct algebraic generating function}
\label{sec:input}
Let
\[
  G(z)=\sum_{n\ge 0}g_nz^n
\]
be the algebraic generating function constructed in \cite{Scheuerle2026}.
It satisfies
\[
  A(z)G(z)^2+B(z)G(z)+C(z)=0,
\]
where the polynomials $A,B,C$ depend explicitly on
$(\lambda,m,r,\eta,\taupar)$. Its associated Hankel transform is initialized by
\[
  H_1=\lambda,\qquad H_2=m,\qquad H_3=r,\qquad H_4=\eta,
\]
and satisfies
\begin{equation}
  H_{n+4}H_n=H_{n+3}H_{n+1}+\taupar H_{n+2}^2.
  \label{eq:somos4}
\end{equation}
Throughout the main workflow we assume that all displayed denominators and
sigma factors are nonzero and that the elliptic discriminant is nonzero.
Exceptional cases are interpreted on the projective curve or by continuation
in the parameters.

The generating-function discriminant factors as
\[
 B(z)^2-4A(z)C(z)
 = (\lambda^2\eta m r^2z)^2
   \Bigl((1-Jz+\taupar z^2)^2-4z^3\Bigr),
\]
where
\begin{equation}
  J=
  \frac{\lambda\eta}{mr}
  +\frac{m^2}{\lambda r}
  +\frac{r^2}{m\eta}
  +\frac{\taupar mr}{\lambda\eta}.
  \label{eq:J}
\end{equation}
Removing the square prefactor and applying the reciprocal change of variable
suggested by the residual polynomial gives the quartic model
\begin{equation}
  \EJ:\qquad y^2=(X^2-JX+\taupar)^2-4X.
  \label{eq:quartic}
\end{equation}
The subsequent workflow begins with \cref{eq:J,eq:quartic}. The explicit
factorization involving $A(z),B(z),C(z)$ belongs to the construction in
\cite{Scheuerle2026}; it is used here as input and is not recomputed by the
appendix script.

\section{What information is already contained in the discriminant?}
\label{sec:discriminant}
The factorization above provides the conceptual bridge between the algebraic
and transcendental descriptions. Once the irrelevant square factor is
removed, the generating-function discriminant supplies a quartic genus-one
model depending only on $(J,\taupar)$. More precisely, $\taupar$ is the
recurrence parameter and the pair $(J,\taupar)$ determines the invariant
curve. The remaining seed information selects a point on that curve, while
the multiplicative gauge normalization enters the constants $A$ and $B$ of
the sigma representation. Different five-tuples may therefore determine the
same pair $(J,\taupar)$ but select different points or gauges.

This observation should be interpreted carefully. At this stage the
discriminant supplies a genus-one model. Its identification with the QRT
invariant curve of the recurrence is established in Workflow Step~4. Even
then, the curve alone does not fix every choice needed for a numerical theta
representation. One must still:
\begin{enumerate}[leftmargin=2.3em]
  \item choose and verify a birational map to a Weierstrass model;
  \item adopt an ordering and orientation for the periods;
  \item identify the point implementing one Somos step;
  \item choose Abel--Jacobi lifts of the initial point and translation;
  \item determine the elementary exponential factors; and
  \item check the result against the original recurrence.
\end{enumerate}
In this precise sense, the theta function is hidden in the discriminant, but
extracting it requires a controlled sequence of normalizations.

\section{The constructive workflow}
\label{sec:workflow}

\workflowstep{1: Compute the invariant and construct the quartic}
Given $(\lambda,m,r,\eta,\taupar)$, compute $J$ using \cref{eq:J} and form the
quartic \cref{eq:quartic}. Before proceeding, test that the quartic is
nonsingular. Equivalently, its polynomial discriminant must be nonzero.
Degenerate inputs require a separate rational or singular analysis and are
outside the principal workflow.

\paragraph{Checkpoint.}
Expand the two sides of the generating-function discriminant identity and
verify exact equality in the coefficient field.

\workflowstep{2: Convert the quartic to Weierstrass form}
Choose as origin the point at infinity characterized by $y/X^2\to1$. Put
\begin{equation}
  U=\frac{y+X^2-JX+\taupar}{2},
  \qquad
  V=U(2X-J)-1.
  \label{eq:quartic-to-cubic}
\end{equation}
Since the quartic equation implies
\[
  U\bigl(U-(X^2-JX+\taupar)\bigr)=-X,
\]
a direct elimination of $X$ and $y$ yields
\begin{equation}
  V^2=4U^3+(J^2-4\taupar)U^2+2JU+1.
  \label{eq:cubic-with-quadratic-term}
\end{equation}
The shift
\begin{equation}
  u=U+\frac{J^2-4\taupar}{12},
  \qquad v=V
  \label{eq:cubic-shift}
\end{equation}
therefore gives the Weierstrass model
\begin{equation}
  v^2=4u^3-g_2u-g_3,
  \label{eq:weierstrass}
\end{equation}
where, with $a=J^2-4\taupar$,
\begin{equation}
  g_2=\frac{a^2}{12}-2J,
  \qquad
  g_3=-\frac{a^3}{216}+\frac{aJ}{6}-1.
  \label{eq:g2g3}
\end{equation}

The inverse affine formulas are equally simple. First set
\[
  U=u-\frac{J^2-4\taupar}{12},\qquad V=v.
\]
For $U\ne0$,
\begin{equation}
  X=\frac{V+JU+1}{2U},
  \qquad
  y=2U-X^2+JX-\taupar.
  \label{eq:cubic-to-quartic}
\end{equation}
Direct symbolic substitution, followed where appropriate by reduction modulo
the quartic or cubic equation, verifies both compositions identically on
their affine charts.
The denominator $U$ vanishes at the two cubic points $(U,V)=(0,1)$ and
$(0,-1)$. In the projective correspondence these are two points omitted by
the affine inverse chart. The chosen origin $y/X^2\to1$ maps to the cubic
point at infinity. The other quartic point at infinity maps to $(U,V)=(0,1)$,
while $(U,V)=(0,-1)$ maps to the finite quartic point
$(X,y)=(0,-\taupar)$. Thus the sole affine exceptional set of the inverse
formula is $U=0$; it is resolved by using the projective models.

The polynomial discriminant of the quartic in $X$ is
\begin{equation}
  \operatorname{disc}_X\bigl((X^2-JX+\taupar)^2-4X\bigr)
  =256\Delta,
  \label{eq:quartic-discriminant}
\end{equation}
where
\begin{equation}
  \Delta=g_2^3-27g_3^2
  =J^4\taupar+J^3-8J^2\taupar^2-36J\taupar
   +16\taupar^3-27.
  \label{eq:elliptic-discriminant}
\end{equation}
Consequently, over a field of characteristic zero, the birational elliptic
model is nonsingular precisely when $\Delta\ne0$. This statement concerns the
curve itself. The computational workflow has additional open conditions from
the recurrence denominators, the affine chart $U\ne0$, and the sigma-function
denominators.

\paragraph{Checkpoint.}
Reduce
\[
  V^2-\bigl(4U^3+(J^2-4\taupar)U^2+2JU+1\bigr)
\]
modulo $y^2-(X^2-JX+\taupar)^2+4X$ and verify that the remainder is zero.
Verify the two inverse compositions, \cref{eq:quartic-discriminant}, and
\[
  j=1728\frac{g_2^3}{\Delta}.
\]

\workflowstep{3: Compute and normalize the period lattice}
Let $e_1,e_2,e_3$ be the distinct roots of
\[
  4u^3-g_2u-g_3=0.
\]
They are the branch values of the double covering
\[
  v^2=4(u-e_1)(u-e_2)(u-e_3).
\]
A period lattice is obtained by integrating the holomorphic differential
$du/v$ over an oriented homology basis:
\[
  2\omega_1=\oint_{\gamma_1}\frac{du}{v},
  \qquad
  2\omega_2=\oint_{\gamma_2}\frac{du}{v}.
\]
The cycles are oriented and, if necessary, replaced by an equivalent basis so
that
\[
  \operatorname{Im}\left(\frac{\omega_2}{\omega_1}\right)>0.
\]
We then set $\taumod=\omega_2/\omega_1$. In this paper, ``normalized'' means
only that the ordered half-period basis has this orientation condition; it
does not, unless stated separately, mean reduction to the standard modular
fundamental domain. Only the lattice is intrinsic; the ordered pair
$(2\omega_1,2\omega_2)$ depends on the chosen homology basis. This choice must
therefore be retained unchanged in the later definitions of elliptic
logarithms, sigma functions, and theta functions.

For real input data with $\Delta>0$, the three roots are real and may be
ordered as $e_1>e_2>e_3$. Define
\[
  k^2=\frac{e_2-e_3}{e_1-e_3},
  \qquad
  k'^2=1-k^2=\frac{e_1-e_2}{e_1-e_3}.
\]
With the positive real square root of $e_1-e_3$, a convenient rectangular
half-period basis is
\begin{equation}
  \omega_1=\frac{K(k)}{\sqrt{e_1-e_3}},
  \qquad
  \omega_2=\frac{iK(k')}{\sqrt{e_1-e_3}},
  \label{eq:periods-three-real-roots}
\end{equation}
where
\[
  K(k)=\int_0^{\pi/2}\frac{d\phi}{\sqrt{1-k^2\sin^2\phi}}.
\]
Consequently, $\taumod=iK(k')/K(k)$ lies in the upper half-plane.

For general complex parameters, or for real parameters with only one real
root, it is preferable to compute the periods from integrals over a chosen
pair of branch cuts rather than to impose the ordering used in
\cref{eq:periods-three-real-roots}. The resulting basis is then normalized by
an element of $\mathrm{SL}_2(\mathbb Z)$ so that
$\operatorname{Im}(\taumod)>0$. A direct elliptic-curve period routine provides an independent numerical check
\cite{PARIGPManual,DLMF23}. In the companion code, PARI's output
$W=(W_1,W_2)=\operatorname{ellperiods}(E)$ is a full-period basis for the same
lattice, whereas the text uses half-periods. For the displayed convention the
code sets $\omega_1=W_1/2$ and $\omega_2=-W_2/2$ and checks
$\operatorname{Im}(\omega_2/\omega_1)>0$.

\paragraph{Checkpoint.}
Verify
\[
  4(u-e_1)(u-e_2)(u-e_3)=4u^3-g_2u-g_3,
  \qquad \operatorname{Im}(\taumod)>0.
\]
Then reconstruct the invariants from the normalized lattice and check that
they agree with the original $g_2$ and $g_3$ to the working precision. In the
three-real-root case, also compare \cref{eq:periods-three-real-roots} with
direct numerical integration of $du/v$.

\workflowstep{4: Determine the elliptic translation point}
Introduce the gauge-invariant recurrence variables
\begin{equation}
  x_n=\frac{H_{n-1}H_{n+1}}{H_n^2}.
  \label{eq:qrt-variable}
\end{equation}
The Somos--4 recurrence \cref{eq:somos4} implies the second-order QRT map
\begin{equation}
  x_{n+2}x_nx_{n+1}^2=x_{n+1}+\taupar,
  \qquad
  x_{n+2}=\frac{x_{n+1}+\taupar}{x_nx_{n+1}^2},
  \label{eq:qrt-map}
\end{equation}
and the quantity
\begin{equation}
  x_nx_{n+1}+\frac{1}{x_n}+\frac{1}{x_{n+1}}
  +\frac{\taupar}{x_nx_{n+1}}
  \label{eq:qrt-invariant}
\end{equation}
is independent of $n$. Its value is precisely $J$ from \cref{eq:J}. Thus the
parameter obtained from the generating-function discriminant is the first
integral of the reduced Somos--4 QRT map, not merely a parameter of an
isomorphic elliptic curve. Indeed,
\[
  x_2=\frac{\lambda r}{m^2},
  \qquad
  x_3=\frac{m\eta}{r^2},
\]
and substituting these values into \cref{eq:qrt-invariant} gives
\cref{eq:J}.

The invariant biquadratic becomes the quartic \cref{eq:quartic} under
\begin{equation}
  X_n=x_nx_{n+1},
  \qquad
  y_n=x_n-x_{n+1}.
  \label{eq:qrt-to-quartic}
\end{equation}
The sign convention in \cref{eq:qrt-to-quartic} is essential: replacing
$y_n$ by $-y_n$ reverses the direction of the translation point below.
Applying \cref{eq:quartic-to-cubic,eq:cubic-shift} gives
\begin{align}
  U_n&=\frac{y_n+X_n^2-JX_n+\taupar}{2},
  &V_n&=U_n(2X_n-J)-1,\nonumber\\
  u_n&=U_n+\frac{J^2-4\taupar}{12},
  &Q_n&=(u_n,V_n)\in E_J.
  \label{eq:qrt-state-on-weierstrass}
\end{align}
With the group law on \cref{eq:weierstrass} and the point at infinity as
identity, one QRT step is
\begin{equation}
  Q_{n+1}=Q_n+P,
  \qquad
  P=\left(\frac{J^2-4\taupar}{12},1\right).
  \label{eq:translation-point}
\end{equation}
Thus the translation point is determined directly by $(J,\taupar)$ and need
not be recovered numerically from several iterates. Here is a compact group-law
derivation. Write $a=x_n$, $b=x_{n+1}$, and
$c=(b+\taupar)/(ab^2)$. On the intermediate cubic of
\cref{eq:cubic-with-quadratic-term}, the state $(a,b)$ has
\[
  (U,V)=(-b,b^2(c-a)),
\]
while the point $P$ has coordinates $(0,1)$. The line through these points has
slope
\[
  M=\frac{V-1}{U}
   =\frac{a^2b^2+a-b-\taupar}{ab}.
\]
Substitution in the chord-and-tangent law for
$V^2=4U^3+(J^2-4\taupar)U^2+2JU+1$ gives the third intersection, after
reflection, as
\[
  U'=-c,\qquad V'=c^2(d-b),
  \qquad d=\frac{c+\taupar}{bc^2}.
\]
These are precisely the coordinates attached to the next QRT state $(b,c)$.
The subsequent shift from $U$ to $u$ leaves the group law unchanged, proving
$Q_{n+1}=Q_n+P$ on the open chart. A direct computational check constructs
$Q_{n+1}$ from
\cref{eq:qrt-map,eq:qrt-to-quartic,eq:qrt-state-on-weierstrass} and verifies
$Q_{n+1}-Q_n=P$ by the Weierstrass group law. The formula is valid on the open
chart where the displayed expressions are defined; exceptional or vanishing
recurrence values are handled projectively or by continuation.

There is also a continued-fraction interpretation. On the quartic model, the
two points at infinity differ by a fixed divisor class. Advancing one line of
the genus-one continued fraction translates the Jacobian by that class
\cite{Hone2020}. With the origin chosen in Workflow Step~2, the second point at
infinity maps to $(U,V)=(0,1)$, hence to $P$ in
\cref{eq:translation-point}. The QRT and continued-fraction routes therefore give the same increment. In
this way the continued-fraction/Hankel construction supplies a complementary
structural explanation for why the discriminant curve and the recurrence
dynamics meet on the same Jacobian.

For the classical data $H_1=H_2=H_3=H_4=1$ and $\taupar=1$, one has $J=4$ and
\[
  E_J:\quad v^2=4u^3-4u+1,
  \qquad P=(1,1).
\]
Furthermore, $x_2=x_3=1$ and \cref{eq:qrt-map} gives $x_4=2$. The
corresponding points are
\[
  Q_2=(0,1),
  \qquad Q_3=(-1,-1),
\]
and the group law gives $Q_3-Q_2=(1,1)=P$.

\paragraph{Checkpoint.}
For several consecutive pairs $(x_n,x_{n+1})$, verify
\cref{eq:qrt-invariant}, map the states to $Q_n$, and check
$Q_{n+1}-Q_n=P$. As a quick structural test, confirm directly that $P$ lies
on \cref{eq:weierstrass}; in the intermediate coordinates it is
$(U,V)=(0,1)$.

\workflowstep{5: Compute Abel--Jacobi coordinates}
Choose lifts $z_2$ and $\delta$ in $\mathbb C$ satisfying
\[
  (\wp(z_2),\wp'(z_2))=Q_2,
  \qquad
  (\wp(\delta),\wp'(\delta))=P,
\]
or, equivalently,
\[
  z_2=\int_{\infty}^{Q_2}\frac{du}{v},
  \qquad
  \delta=\int_{\infty}^{P}\frac{du}{v}\pmod{\Lambda}.
\]
Then set
\begin{equation}
  z_0=z_2-3\delta \pmod{\Lambda}.
  \label{eq:z0-from-z2}
\end{equation}
The lifts are not unique. Their representatives should be reduced to a
stated fundamental parallelogram, with signs fixed by the $v$-coordinates.
With the normalization used in the companion script, the geometric point
indexed by $Q_n$ is represented by
\[
  z_n=z_0+(n+1)\delta,
\]
so that $Q_{n+1}-Q_n=P$ is encoded by the increment $\delta$ in the Abel--Jacobi
parameter.

\paragraph{Checkpoint.}
For several integers $n$, verify that
\[
  (\wp(z_0+(n+1)\delta),\wp'(z_0+(n+1)\delta))
\]
agrees with the Weierstrass image of $Q_n$.

\workflowstep{6: Fix the sigma-function normalization}
In the convention used here, the classical sigma-function solution of a
nondegenerate Somos--4 recurrence is
\begin{equation}
  H_n=A B^n
  \frac{\sigma(z_0+n\delta)}{\sigma(\delta)^{n^2}}.
  \label{eq:sigma}
\end{equation}
This is the normalization of the general sigma representation proved in
\cite{Hone2005}. The sign of $\delta$ is fixed in Workflow Steps~4 and~5 by
requiring $Q_{n+1}=Q_n+P$ and $(\wp(\delta),\wp'(\delta))=P$. With the
convention used here, the Abel--Jacobi coordinate of $Q_n$ is
\begin{equation}
  z_n=z_0+(n+1)\delta,
  \qquad
  (\wp(z_n),\wp'(z_n))=Q_n.
  \label{eq:sigma-indexing}
\end{equation}
In particular, the initial recurrence data determine $Q_2$ from $x_2$ and
$x_3$, so $z_0$ is obtained from \cref{eq:z0-from-z2}.

\begin{quote}
\textbf{Indexing warning.} The point $Q_n$ associated with the QRT state
$(x_n,x_{n+1})$ corresponds to $z_0+(n+1)\delta$, whereas the sigma
representation of $H_n$ contains $z_0+n\delta$. Consequently
$z_0=z_2-3\delta$ for the present $H_1$-based indexing. Using $z_2$ itself as
$z_0$ would displace the sigma formula by three recurrence indices.
\end{quote}

For brevity, put
\begin{equation}
  S_n=\frac{\sigma(z_0+n\delta)}{\sigma(\delta)^{n^2}}.
  \label{eq:Sn}
\end{equation}
Assuming the displayed sigma factors are nonzero, $H_1=\lambda$ and $H_2=m$
determine the gauge constants uniquely:
\begin{equation}
  B=\frac{mS_1}{\lambda S_2},
  \qquad
  A=\frac{\lambda^2S_2}{mS_1^2}.
  \label{eq:AB-from-S}
\end{equation}
Equivalently,
\begin{align}
  B&=\frac{m\,\sigma(z_0+\delta)\sigma(\delta)^3}
          {\lambda\,\sigma(z_0+2\delta)},
  \label{eq:B-explicit}\\
  A&=\frac{\lambda^2\sigma(z_0+2\delta)}
          {m\,\sigma(\delta)^2\sigma(z_0+\delta)^2}.
  \label{eq:A-explicit}
\end{align}
The remaining initial values are consistency conditions:
\begin{equation}
  r=A B^3S_3,
  \qquad
  \eta=A B^4S_4.
  \label{eq:sigma-initial-checks}
\end{equation}
Thus all four initial values are used in a four-data consistency test: the
first two fix $A$ and $B$, while the third and fourth independently test the
elliptic lift, geometric data, and branch choices.

The exponent $n^2$ in \cref{eq:sigma} is fixed by the sigma addition formula
\begin{equation}
  \frac{\sigma(z+\delta)\sigma(z-\delta)}
       {\sigma(z)^2\sigma(\delta)^2}
  =\wp(\delta)-\wp(z).
  \label{eq:sigma-addition}
\end{equation}
Substitution of \cref{eq:sigma} into \cref{eq:qrt-variable} gives
\begin{equation}
  x_n=
  \frac{\sigma(z_0+(n-1)\delta)\sigma(z_0+(n+1)\delta)}
       {\sigma(z_0+n\delta)^2\sigma(\delta)^2}
  =\wp(\delta)-\wp(z_0+n\delta).
  \label{eq:xn-sigma}
\end{equation}
This identity links the sigma normalization directly to the QRT variables.

\paragraph{Checkpoint.}
First verify \cref{eq:sigma-initial-checks}. Then compare
\cref{eq:sigma} with exact recurrence values for a range of positive and,
where defined, negative indices. Independently check \cref{eq:xn-sigma}; this
detects an incorrect sign of $\delta$, an index shift in $z_0$, or an
inconsistent period representative.

\workflowstep{7: Convert the sigma formula to Jacobi theta form}
For the chosen ordered half-periods, define
\[
  \taumod=\frac{\omega_2}{\omega_1},
  \qquad
  \eta_1=\zeta(\omega_1),
  \qquad
  \vartheta'_1=\theta_1'(0\mid\taumod),
  \qquad
  c=\frac{2\omega_1}{\pi\vartheta'_1}.
\]
The classical sigma--theta identity is
\begin{equation}
 \sigma(z)=
 \exp\!\left(\frac{\eta_1z^2}{2\omega_1}\right)
 c\,
 \theta_1\!\left(\frac{\pi z}{2\omega_1}\,\middle|\,\taumod\right),
 \label{eq:sigma-theta}
\end{equation}
with the standard analytic normalization of $\theta_1$; see
\cite{WhittakerWatson,DLMF23}. Put
\begin{equation}
  A_0=\frac{\pi\delta}{2\omega_1},
  \qquad
  B_0=\frac{\pi z_0}{2\omega_1}.
  \label{eq:theta-arguments}
\end{equation}
Substituting \cref{eq:sigma-theta} into \cref{eq:sigma} and collecting the
constant, linear, and quadratic powers of $n$ gives the explicit theta form
\begin{equation}
  H_n=K_0q_0^{n^2}r_0^n
  \theta_1(A_0n+B_0\mid\taumod),
  \label{eq:theta}
\end{equation}
where
\begin{align}
  K_0&=Ac\exp\!\left(\frac{\eta_1z_0^2}{2\omega_1}\right),
  \label{eq:K0}\\
  q_0&=\exp\!\left(\frac{\eta_1\delta^2}{2\omega_1}\right)\sigma(\delta)^{-1}
       =\bigl(c\,\theta_1(A_0\mid\taumod)\bigr)^{-1},
  \label{eq:q0}\\
  r_0&=B\exp\!\left(\frac{\eta_1z_0\delta}{\omega_1}\right).
  \label{eq:r0}
\end{align}
The exponential in \cref{eq:q0} is essential. Equivalently, applying
\cref{eq:sigma-theta} at $z=\delta$ gives the second expression for $q_0$.

\paragraph{Checkpoint.}
Evaluate \cref{eq:sigma} and \cref{eq:theta} independently and compare them
before comparing either expression with the recurrence. Software libraries
often use different argument scalings or nomes for Jacobi theta functions;
therefore \cref{eq:sigma-theta,eq:theta-arguments} should be treated as the
normalization contract for any implementation.

\section{Algorithm summary}
\begin{enumerate}[label=\textbf{Step \arabic*.},leftmargin=2.7em]
  \item Read $(\lambda,m,r,\eta,\taupar)$ and test the required nonvanishing
        assumptions.
  \item Compute $J$ from \cref{eq:J}.
  \item Construct $\EJ$ and test nonsingularity.
  \item Compute the birational Weierstrass model and retain both maps.
  \item Compute $(g_2,g_3)$ and check curve invariants.
  \item Compute $(\omega_1,\omega_2)$ and normalize the period basis.
  \item Construct consecutive QRT states and verify the fixed point $P$.
  \item Compute and normalize $z_0$ and $\delta$.
  \item Determine $A$ and $B$ from the initial data and check $H_3,H_4$.
  \item Convert the result to theta form using
        \cref{eq:sigma-theta,eq:theta-arguments,eq:K0,eq:q0,eq:r0}.
  \item Validate recurrence, generating-function discriminant, sigma form,
        and theta form.
\end{enumerate}

\section{Numerical examples}
\label{sec:examples}

\subsection{The classical Somos--4 sequence}
For
\[
  (\lambda,m,r,\eta,\taupar)=(1,1,1,1,1),
\]
the companion script computes
\[
  J=4,
  \qquad
  (g_2,g_3,\Delta)=(4,-1,37).
\]
The corresponding quartic is
\[
  y^2=(X^2-4X+1)^2-4X.
\]
One compatible pair of half-periods, obtained at 80 decimal digits of working
precision and oriented so that $\operatorname{Im}(\omega_2/\omega_1)>0$, is
\[
  \omega_1=1.4967293231159798149,
  \qquad
  \omega_2=1.2256946909933950304\,i.
\]
The corresponding Abel--Jacobi parameters are
\[
  z_0=-2.7887781458561870232+1.2256946909933950304\,i,
  \qquad
  \delta=1.8591854305707913488.
\]
The sigma-normalization constants are
\[
  A=-0.23674493915977148+0.06797467161582982\,i,
\]
\[
  B=1.24599348339577469+3.55401276446268772\,i.
\]
Although these constants and lifts are complex, their combination in
\cref{eq:sigma} gives real integral values to numerical precision. The complex
phases arise from the chosen representatives of $z_0$ and $\delta$ and from
the selected period basis; consistent changes of these choices leave the
resulting sequence unchanged. The result
is
\[
  H_1,H_2,\ldots=1,1,1,1,2,3,7,23,59,\ldots,
\]
the classical Somos--4 sequence, OEIS A006720 \cite{OEISA006720}.

\subsection{A nontrivial member of the five-parameter family}
For $(\lambda,m,r,\eta,\taupar)=(2,1,1,3,2)$, exact arithmetic gives
\[
  J=\frac{43}{6},
  \qquad
  (g_2,g_3,\Delta)=
  \left(\frac{2213809}{15552},
  -\frac{3291850729}{10077696},
  \frac{1161689}{324}\right).
\]
The first terms of the orbit are
\[
  2,1,1,3,\frac{5}{2},\frac{41}{2},74,\frac{2051}{6},
  \frac{43103}{6},\frac{335987}{9},\ldots.
\]
This nonunit rational seed shows that the exact algebraic part of the workflow
remains consistent away from the symmetric seed $(1,1,1,1,1)$ before the
numerical period and sigma evaluation is performed.

\section{Computational reliability and reproducibility}
The workflow mixes exact algebra and high-precision complex analysis. Its
exceptional cases fall into distinct classes: $\Delta=0$ makes the curve
singular; vanishing recurrence values can make a QRT or recurrence denominator
undefined; $U=0$ leaves the affine inverse chart and must be treated
projectively; and a vanishing sigma denominator requires a different
normalization or continuation. The main elliptic workflow excludes these
situations. Subject to those conditions, the following practices are
recommended:
\begin{itemize}[leftmargin=2em]
  \item perform discriminant factorization and birational-map checks exactly
        whenever the input data are rational;
  \item postpone floating-point conversion until period computation;
  \item state whether reported periods are full periods or half-periods;
  \item state the orientation and period-basis convention;
  \item reduce Abel--Jacobi arguments consistently modulo the same lattice;
  \item test signs with the second Weierstrass coordinate, not only the
        $u$-coordinate;
  \item when an independent Jacobi-theta implementation is available, compare
        the theta conversion with the sigma representation; and
  \item validate several consecutive and nonconsecutive indices.
\end{itemize}
These checks are not cosmetic. A wrong sign of $\delta$, a shifted lift of
$z_0$, or a modularly equivalent but inconsistently transformed period basis
can produce formulas that appear plausible for a few terms while failing
globally.

\section{Discussion and future directions}
The discriminant of the direct algebraic generating function provides a
particularly transparent entry point into the elliptic geometry of a
Somos--4 orbit. The present workflow is intended to make that geometry
computationally accessible rather than to re-establish the classical
elliptic-function theory. It also clarifies which parts of the transcendental
representation are intrinsic to the curve and which record the initial
point, indexing, and normalization.

Natural extensions include:
\begin{itemize}[leftmargin=2em]
  \item direct extraction of the translation point from paired-Hankel or
        continued-fraction data;
  \item systematic formulas for modular invariants in terms of
        $(J,\taupar)$;
  \item arithmetic analysis of special parameter choices;
  \item analogous workflows for Somos--5 and Somos--6; and
  \item higher-genus generalizations involving Kleinian sigma functions.
\end{itemize}

\section{Conclusion}
The elliptic curve needed for a sigma- or theta-function description is
already present in the discriminant of the direct algebraic generating
function. What remains is a sequence of explicit, verifiable conversions. By
collecting those conversions, fixing their conventions, and implementing
them in one script, the scattered classical theory becomes a usable tool for
concrete Somos--4 sequences.

\appendix
\section{PARI/GP companion implementation}
\label{app:pari}
The code below carries the normalization used in the text through the exact
algebraic checks beginning with $(J,\taupar)$, the period calculation, the QRT
translation check, and the elliptic sigma evaluation. It does not reproduce
the original $A(z),B(z),C(z)$ discriminant factorization, which is taken from
\cite{Scheuerle2026}, and it does not independently evaluate the Jacobi theta
function. PARI/GP represents
\[
  v^2=4u^3-g_2u-g_3
\]
by
\[
  Y^2=u^3-\frac{g_2}{4}u-\frac{g_3}{4},
  \qquad Y=\frac{v}{2}.
\]
Thus a text point $(u,v)$ is supplied to PARI/GP as $[u,v/2]$. The functions
used below are documented in \cite{PARIGPManual}.

\begin{lstlisting}[style=parigp]
\\ PARI/GP companion for the algebraic-to-sigma workflow.

somos_J(lambda, m, r, eta, tau) =
{
  lambda*eta/(m*r)
  + m^2/(lambda*r)
  + r^2/(m*eta)
  + tau*m*r/(lambda*eta);
};

quartic_poly(J, tau, X) = (X^2 - J*X + tau)^2 - 4*X;

qrt_invariant(x, y, tau) = x*y + 1/x + 1/y + tau/(x*y);
qrt_step(x, y, tau) = [y, (y + tau)/(x*y^2)];
qrt_step_pair(x, y, tau) = qrt_step(x, y, tau);

somos_recurrence_step(v, tau) =
{
  my(n = #v);
  if (n < 4, error("Need at least four values to advance the Somos-4 orbit"));
  if (v[n-3] == 0, error("Zero denominator in Somos-4 recurrence"));
  (v[n]*v[n-2] + tau*v[n-1]^2) / v[n-3];
};

somos_orbit(lambda, m, r, eta, tau, N) =
{
  my(v = List([lambda, m, r, eta]));
  if (N < 1, return([]));
  for (n = 5, N,
    listput(v, somos_recurrence_step(v, tau));
  );
  Vec(v)[1..min(N, #v)];
};

quartic_to_weierstrass(X, Y, J, tau) =
{
  my(U = (Y + X^2 - J*X + tau)/2);
  my(V = U*(2*X - J) - 1);
  [U + (J^2 - 4*tau)/12, V];
};

weierstrass_to_quartic(u, v, J, tau) =
{
  my(U = u - (J^2 - 4*tau)/12);
  if (U == 0, error("Affine inverse map hits the excluded U = 0 chart"));
  my(X = (v + J*U + 1)/(2*U));
  [X, 2*U - X^2 + J*X - tau];
};

weierstrass_invariants(J, tau) =
{
  my(a = J^2 - 4*tau);
  my(g2 = a^2/12 - 2*J);
  my(g3 = -a^3/216 + a*J/6 - 1);
  [g2, g3, g2^3 - 27*g3^2];
};

ell_curve_from_J(J, tau) =
{
  my(w = weierstrass_invariants(J, tau));
  ellinit([0, 0, 0, -w[1]/4, -w[2]/4]);
};

fixed_translation_point(J, tau) =
{
  my(a = J^2 - 4*tau);
  [a/12, 1/2];       \\ PARI coordinates [u,v/2]
};

translation_point_on_curve(J, tau) = fixed_translation_point(J, tau);

curve_point_from_qrt(x, y, J, tau) =
{
  my(P = quartic_to_weierstrass(x*y, x-y, J, tau));
  [P[1], P[2]/2];    \\ PARI coordinates [u,v/2]
};

check_discriminant(J, tau) =
{
  my(w = weierstrass_invariants(J, tau));
  my(F = quartic_poly(J, tau, 'X));
  [poldisc(F, 'X), 256*w[3], poldisc(F, 'X) - 256*w[3]];
};

check_quartic_weierstrass(J, tau) =
{
  my(w = weierstrass_invariants(J, tau));
  my(E = ell_curve_from_J(J, tau));
  my(W = ellperiods(E));
  my(omega1 = W[1]/2, omega2 = -W[2]/2);
  my(tau_mod = omega2/omega1);

  print("J = ", J, "; g2 = ", w[1], "; g3 = ", w[2], "; Delta = ", w[3]);
  print("PARI full periods W = ", W);
  print("Reference half-periods: omega1 = ", omega1, ", omega2 = ", omega2);
  print("tau_mod = ", tau_mod);
  print("discriminant check = ", check_discriminant(J, tau));

  if (imag(tau_mod) <= 0,
    print("Warning: tau_mod is not in the upper half-plane.");
  );

  abs(check_discriminant(J, tau)[3]) < 1e-30;
};

verify_qrt_translation(lambda, m, r, eta, tau, nsteps = 6) =
{
  my(J = somos_J(lambda, m, r, eta, tau));
  my(E = ell_curve_from_J(J, tau));
  my(x = lambda*r/m^2, y = m*eta/r^2);
  my(P = fixed_translation_point(J, tau));
  for (k = 1, nsteps,
    my(Qn = curve_point_from_qrt(x, y, J, tau));
    my(next = qrt_step(x, y, tau));
    my(Qnp1 = curve_point_from_qrt(next[1], next[2], J, tau));
    print("step ", k,
          ": invariant error = ", qrt_invariant(x, y, tau) - J,
          ", translation exact = ", ellsub(E, Qnp1, Qn) == P);
    x = next[1];
    y = next[2];
  );
};

sigma_lattice_from_J(J, tau) =
{
  my(E = ell_curve_from_J(J, tau));
  ellperiods(E);
};

sigma_representation(lambda, m, r, eta, tau, z0, delta) =
{
  my(L = sigma_lattice_from_J(somos_J(lambda, m, r, eta, tau), tau));
  my(sd = ellsigma(L, delta));
  my(S1 = ellsigma(L, z0 + delta));
  my(S2 = ellsigma(L, z0 + 2*delta));
  my(B = m*S1*sd^3 / (lambda*S2));
  my(A = lambda^2*S2 / (m*sd^2*S1^2));
  [A, B, sd, S1, S2];
};

sigma_parameters(lambda, m, r, eta, tau) =
{
  my(J = somos_J(lambda, m, r, eta, tau));
  my(E = ell_curve_from_J(J, tau));
  my(x2 = lambda*r/m^2, x3 = m*eta/r^2);
  my(Q2 = curve_point_from_qrt(x2, x3, J, tau));
  my(P = fixed_translation_point(J, tau));
  my(L = ellperiods(E));
  my(delta = ellpointtoz(E, P));
  my(z2 = ellpointtoz(E, Q2));
  \\ In the H_1-started normalization, Q_n corresponds to z0 + (n+1)*delta.
  my(z0 = z2 - 3*delta);
  my(sd = ellsigma(L, delta));
  my(S1 = ellsigma(L, z0 + delta)/sd);
  my(S2 = ellsigma(L, z0 + 2*delta)/sd^4);
  my(B = m*S1/(lambda*S2));
  my(A = lambda^2*S2/(m*S1^2));
  [E, L, z0, delta, A, B];
};

sigma_value(sigpar, n) =
{
  my(L = sigpar[2], z0 = sigpar[3], delta = sigpar[4]);
  my(A = sigpar[5], B = sigpar[6]);
  A*B^n*ellsigma(L, z0 + n*delta)/ellsigma(L, delta)^(n^2);
};

somos_sigma_value(lambda, m, r, eta, tau, z0, delta, n) =
{
  my(L = sigma_lattice_from_J(somos_J(lambda, m, r, eta, tau), tau));
  my(A, B, sd, S1, S2);
  [A, B, sd, S1, S2] = sigma_representation(lambda, m, r, eta, tau, z0, delta);
  A*B^n*ellsigma(L, z0 + n*delta)/sd^(n^2);
};

workflow_demo(lambda, m, r, eta, tau, N = 10, prec = 80) =
{
  default(realprecision, prec);
  my(J = somos_J(lambda, m, r, eta, tau));
  my(w = weierstrass_invariants(J, tau));
  if (w[3] == 0, error("Singular quartic / elliptic curve"));
  my(sigpar = sigma_parameters(lambda, m, r, eta, tau));
  my(exact = somos_orbit(lambda, m, r, eta, tau, N));
  my(approx = vector(N, n, sigma_value(sigpar, n)));
  my(err = vector(N, n, abs(approx[n] - exact[n])));
  my(k0 = min(4, N));
  my(initial_expected = [lambda, m, r, eta][1..k0]);
  my(initial_sigma = approx[1..k0]);
  my(initial_errors = vector(k0, n,
      abs(initial_sigma[n] - initial_expected[n]));
  print("J = ", J);
  print("[g2,g3,Delta] = ", w);
  print("ellperiods(E) = ", ellperiods(sigpar[1]));
  print("z0 = ", sigpar[3], "; delta = ", sigpar[4]);
  print("A = ", sigpar[5], "; B = ", sigpar[6]);
  print("exact orbit = ", exact);
  print("sigma orbit = ", approx);
  print("initial four sigma errors = ", initial_errors);
  print("absolute errors = ", err);
  [J, w, ellperiods(sigpar[1]), sigpar, exact, approx, err];
};

demo_sigma_workflow() =
{
  workflow_demo(1, 1, 1, 1, 1, 8, 80);
};

check_sigma_workflow_classical() =
{
  workflow_demo(1, 1, 1, 1, 1, 10, 80);
};

classical_example() =
{
  workflow_demo(1, 1, 1, 1, 1, 10, 80);
};

nonunit_example() =
{
  workflow_demo(2, 1, 1, 3, 2, 10, 80);
};
\end{lstlisting}

The script deliberately performs the recurrence and sigma evaluations through
separate code paths. The displayed error vector is therefore a direct check
of the indexing, elliptic logarithms, period convention, and gauge factors.
The theta representation is obtained algebraically from the sigma identity.
It remains available for an independent numerical check in software that
implements the standard Jacobi $\theta_1(z\mid\tau)$ convention, using
\cref{eq:theta-arguments,eq:K0,eq:q0,eq:r0}; that independent theta evaluation
is not part of the supplied PARI/GP script.

\section{Verification checklist}
\begin{enumerate}[leftmargin=2.4em]
  \item The generating-function discriminant factorization is taken from
        \cite{Scheuerle2026}; it is not recomputed by the appendix script.
  \item The quartic constructed from $(J,\taupar)$ is nonsingular.
  \item Forward and inverse birational maps are verified symbolically.
  \item Curve discriminants and $j$-invariants agree under the stated
        normalization.
  \item The period basis has $\operatorname{Im}(\taumod)>0$.
  \item The computed point difference $Q_{n+1}-Q_n$ is independent of $n$.
  \item The signs of elliptic logarithms agree with the $v$-coordinates.
  \item Sigma values satisfy all four initial conditions and the recurrence.
  \item The Jacobi-theta formula follows algebraically from the sigma formula
        and is available for independent numerical verification.
  \item Reported numerical errors are compatible with the working precision.
\end{enumerate}

\section{Statement on AI assistance and computational tools}
This manuscript was prepared with AI-assisted language editing, \LaTeX{}
formatting, and development of verification code. The author reviewed the
mathematical statements, normalization choices, references, and numerical
outputs. PARI/GP was used for exact algebra and high-precision elliptic
computations; its relevant conventions are documented in
\cite{PARIGPManual}.

\bibliographystyle{plain}
\begin{thebibliography}{99}

\bibitem{Scheuerle2026}
Thomas Scheuerle,
\newblock \emph{Direct Generation of a Somos--4 Sequence from an Algebraic
Generating Function},
\newblock arXiv:2609.15754 [math.CO], 2026.
\newblock \url{https://arxiv.org/abs/2609.15754}.

\bibitem{Hone2005}
A.~N.~W. Hone,
\newblock Elliptic curves and quadratic recurrence sequences,
\newblock \emph{Bulletin of the London Mathematical Society} 37 (2005),
161--171.
\newblock \url{https://doi.org/10.1112/S0024609304004163}.

\bibitem{HoneSwart2008}
A.~N.~W. Hone and C.~S. Swart,
\newblock Integrality and the Laurent phenomenon for Somos 4 and Somos 5
sequences,
\newblock \emph{Mathematical Proceedings of the Cambridge Philosophical
Society} 145 (2008), 65--85.
\newblock \url{https://doi.org/10.1017/S030500410800114X}.

\bibitem{VanderPoortenSwart2006}
A.~J. van der Poorten and C.~S. Swart,
\newblock Recurrence relations for elliptic sequences: every Somos 4 is a
Somos $k$,
\newblock \emph{Bulletin of the London Mathematical Society} 38 (2006),
546--554.
\newblock \url{https://doi.org/10.1112/S0024609306018534}.

\bibitem{Ward1948}
Morgan Ward,
\newblock Memoir on elliptic divisibility sequences,
\newblock \emph{American Journal of Mathematics} 70 (1948), 31--74.

\bibitem{Hone2020}
A.~N.~W. Hone,
\newblock Continued fractions and Hankel determinants from hyperelliptic
curves,
\newblock \emph{Communications on Pure and Applied Mathematics} 73 (2020),
2310--2347.
\newblock \url{https://doi.org/10.1002/cpa.21923}.

\bibitem{WhittakerWatson}
E.~T. Whittaker and G.~N. Watson,
\newblock \emph{A Course of Modern Analysis}, 4th ed.,
\newblock Cambridge University Press, Cambridge, 1927; reprinted 1996.

\bibitem{DLMF23}
NIST Digital Library of Mathematical Functions,
\newblock Chapter 23: Weierstrass elliptic and modular functions,
\newblock \url{https://dlmf.nist.gov/23}, accessed September 16, 2026.

\bibitem{PARIGPManual}
The PARI Group,
\newblock \emph{User's Guide to PARI/GP},
\newblock \url{https://pari.math.u-bordeaux.fr/doc.html}, accessed
September 16, 2026.

\bibitem{OEISA006720}
N.~J.~A. Sloane and The OEIS Foundation Inc.,
\newblock \emph{The On-Line Encyclopedia of Integer Sequences, A006720:
Somos--4 sequence},
\newblock \url{https://oeis.org/A006720}, accessed September 16, 2026.

\end{thebibliography}
\end{document}
